\documentclass[11pt]{article}
\usepackage[a4paper,margin=27mm]{geometry}
\usepackage{amsmath,amssymb,amsthm,hyperref}
\newtheorem{theorem}{Theorem}
\newtheorem{lemma}[theorem]{Lemma}
\newcommand{\Z}{\mathbb Z}
\newcommand{\Q}{\mathbb Q}
\title{A nearly sharp rational quadrature upper bound from Mit'kin}
\author{A consequence of published Hilbert--Kamke and real quadrature results}
\date{30 September 2026}
\begin{document}
\maketitle
Put
\[
 H(r)=\sum_{a=0}^{\lfloor\log_2r\rfloor}
 2^a(2^{\lfloor r/2^a\rfloor}-1)=2^r+O(2^{r/2}).
\]
Let $R_r$ be the minimum size of a rational equal-weight uniform-interval
quadrature rule of degree $r$, allowing repeated nodes.
\begin{theorem}\label{main}
For every sufficiently large $r$, every integer $K\ge H(r)$ admits
rational nodes $x_1,\ldots,x_K\in(0,1)$ with
\[
 K^{-1}\sum_i x_i^j=1/(j+1)\qquad(0\le j\le r).
\]
Thus $R_r\le2^r+O(2^{r/2})$. Together with the independently proved
optimized parity-filter bound $R_r\ge2^{2\lfloor r/2\rfloor}$,
this gives $2^r\le R_r\le2^r+O(2^{r/2})$ for even $r$.
\end{theorem}
The upper bound is derived from earlier published results; no historical
novelty claim is made. The only adaptation is to keep the circle-method
summation interval fixed at $[0,P]$, rather than tying its endpoint to a
prescribed power sum. Unlocalized positive-integer solutions alone
would not ensure rational nodes in $[0,1]$.

\begin{lemma}[Mit'kin's uniform local input]\label{local}
Let $r\ge12$, $K\ge H(r)$, and $V=(h^j)_{1\le j,h\le r}$.
If $N\in\Z^r$ and $Vt=N$ has an integer solution, then the
Hilbert--Kamke singular series $\mathfrak S_{K,r}(N)$ is bounded below
by a positive constant depending only on $r,K$.
\end{lemma}
\begin{proof}[Source and precise extraction]
Mit'kin \cite[Lemma 15, printed p.~577]{mitkin} proves uniform positivity
whenever $K\ge R(r)\ge r^2+2$ and all local lifting thresholds
$\widetilde r(r,p)$ are at most $R(r)$. The estimates assembled on
printed pp.~589--590 below the main theorem give
$\widetilde r(r,p)\le H(r)$ for every prime $p$ when $r\ge12$.
The $p=2$ bound is the corollary to Proposition 2 on p.~589; the
remaining bounds in (94) are at most $H(r)$.
The arithmetic condition is precisely integer solvability of $Vt=N$,
as stated in the introduction and reformulated in (13)--(15).
Also $H(r)\ge r^2+2$ in this range.
\end{proof}

\begin{lemma}[Fixed-interval circle method]\label{circle}
Fix $r\ge3$ and $K\ge r(r+1)+1$. Let $c\in\Q^r$ be fixed, and let
$P$ tend to infinity through integers for which $N_j=c_jP^j$ are
integers. The number $A(P,N)$ of solutions
\[
 1\le y_i\le P-1,\qquad\sum_i y_i^j=N_j\quad(1\le j\le r)
\]
satisfies, with $D=r(r+1)/2$,
\[
 A(P,N)=P^{K-D}\bigl(\mathfrak S_{K,r}(N)\mathfrak J_{K,r}(c)+o(1)\bigr),
\]
where
\[
 \mathfrak J_{K,r}(c)=\int_{\mathbb R^r}
 \left(\int_0^1e(\beta_1t+\cdots+\beta_rt^r)\,dt\right)^K
 e(-\beta\cdot c)\,d\beta.
\]
If the real moment system has an interior solution with $r$ distinct
coordinates, then $\mathfrak J_{K,r}(c)>0$.
\end{lemma}
\begin{proof}
We extract the standard calculation from \cite[Sections 3--6]{wooley},
keeping the endpoint $P$ independent of the targets. Put
\[
 f(\alpha)=\sum_{1\le y\le P-1}e(\alpha_1y+\cdots+\alpha_ry^r).
\]
By orthogonality $A(P,N)=\int_{[0,1)^r}f(\alpha)^K
 e(-\alpha\cdot N)\,d\alpha$. The target factor has modulus one
in all the following estimates.

Take the major arcs $\mathfrak P=\mathfrak K(L)$ of
\cite[Section 3]{wooley}, with $L=P^{1/(8r^2)}$: they have
$q\le L$, $(q,a_1,\ldots,a_r)=1$, and
$|\alpha_j-a_j/q|\le LP^{-j}$ for all $j$.
Lemma 4.1 of that paper gives
\[
 \sup_{\alpha\notin\mathfrak P}|f(\alpha)|\ll P^{1-\delta},
 \qquad\delta=1/(12r^3).
\]
Deleting the two endpoints changes the sum by at most two. The
Vinogradov mean value theorem, as in (1.3) of \cite{wooley}, gives
\[
 \int_{[0,1)^r}|f(\alpha)|^{2D}\,d\alpha\ll_\epsilon P^{D+\epsilon}.
\]
Restricting the summation interval only decreases this even moment,
which counts solutions. Hence the absolute minor-arc contribution is
\[
 \ll P^{K-D-\delta(K-2D)+\epsilon}=o(P^{K-D}),
\]
with $0<\epsilon<\delta(K-2D)$.

On a major arc, the approximation of \cite[Section 6]{wooley} is
\[
 f(\alpha)=q^{-1}S(q,a)I(\alpha-a/q;P)+O(L^2),
\]
where $S(q,a)=\sum_{h=1}^qe((a_1h+\cdots+a_rh^r)/q)$ and
$I(\beta;P)=\int_0^Pe(\beta_1t+\cdots+\beta_rt^r)\,dt$.
The major arcs have total measure $O(L^{2r+1}P^{-D})$, so replacement
of $f^K$ has total error
\[
 O(P^{K-D-1}L^{2r+3})=o(P^{K-D}).
\]
The remaining expression is the truncated singular series times the
truncated singular integral. The estimates used in
\cite[Section 6]{wooley},
\[
 |S(q,a)|\ll_\epsilon q^{1-1/r+\epsilon},\qquad
 |I(\beta;1)|\ll(1+|\beta_1|+\cdots+|\beta_r|)^{-1/r},
\]
give uniformly negligible tails: absolute convergence follows from
$K>r(r+1)$ for the series and $K>r^2$ for the integral. The change
$\beta_j\mapsto P^{-j}\beta_j$ leaves $c_j=N_j/P^j$ fixed, proving
the asymptotic. None of these steps requires $P$ to equal the largest
root of a target moment.

At an interior real solution with $r$ distinct coordinates, the moment
Jacobian has an invertible Vandermonde minor. The implicit function
theorem supplies a positive density at $c$ from a small box about that
solution. Fourier inversion, justified by absolute convergence, identifies
the whole nonnegative density with $\mathfrak J_{K,r}(c)$, proving
strict positivity.
\end{proof}

\begin{proof}[Proof of Theorem~\ref{main}]
Set $c_j=K/(j+1)$, $B=|\det V|$, and restrict $P$ to multiples of
$B(r+1)!$. Then every $N_j=KP^j/(j+1)$ is an integer divisible by $B$,
so $V^{-1}N=(\operatorname{adj}V)N/\det V$ is integral.
Lemma~\ref{local} gives a uniform positive singular-series lower bound.

For sufficiently large $r$, $H(r)$ exceeds both $r(r+1)+1$ and the
quadratic node threshold in \cite{gp} for a real equal-weight rule of
degree $2r+2$. Such a rule has at least $r+2$ distinct support points:
otherwise the square of a polynomial vanishing on its support contradicts
exactness. At least $r$ support points are interior. Fix these as pivots,
move all endpoint coordinates slightly inward, and use the implicit
function theorem on the first $r$ moments to correct the pivots. This
gives an interior real solution with nonsingular Jacobian.
Thus $\mathfrak J_{K,r}(c)>0$.

Lemma~\ref{circle} now gives $A(P,N)>0$ for sufficiently large multiples
$P$. The rational numbers $x_i=y_i/P$ lie in $(0,1)$ and have the required
moments.
\end{proof}

\paragraph{Checks and scope.}
Mit'kin's Lemma 15 and final theorem were checked visually in the
original English PDF, printed pp.~577 and 589. The circle-method lemma
is an explicit deduction from the cited estimates, not a quotation of
Wooley's Theorem 1.1. No common-denominator bound is asserted.
This concerns scalar rational quadrature and does not close the
exponential versus factorial-size gap for permutation-fair dice.

\begin{thebibliography}{9}
\bibitem{mitkin} D.~A.~Mit'kin, \emph{An estimate for the number of terms
in the Hilbert--Kamke problem}, Math. USSR-Sbornik 57 (1987), 561--590,
\href{https://doi.org/10.1070/SM1987v057n02ABEH003087}{doi:10.1070/SM1987v057n02ABEH003087}.
\bibitem{wooley} Trevor D.~Wooley, \emph{Subconvexity and the
Hilbert--Kamke problem}, Math. Z. (2025),
\href{https://doi.org/10.1007/s00209-025-03812-9}{doi:10.1007/s00209-025-03812-9};
\href{https://www.math.purdue.edu/~twooley/publ/20211215hilkam.pdf}{author's full text}.
\bibitem{gp} Shoni Gilboa and Ron Peled, \emph{Chebyshev-type quadratures
for doubling weights}, Constructive Approximation 45 (2017), 193--216,
\href{https://arxiv.org/abs/1507.01505}{arXiv:1507.01505}.
\end{thebibliography}
\end{document}
